\begin{solucion}{Caso base}
Definimos:
\[
\begin{array}{c|l}
\toprule
\text{Simbolo} & \text{Significado} \\
\midrule
q\in\mathbb{R}_+ & \text{cantidad vendida al cliente final.}\\
p\in\mathbb{R}_+ & \text{precio final al cliente.}\\
w\in\mathbb{R}_+ & \text{precio mayorista elegido por el proveedor.}\\
c\in\mathbb{R}_+ & \text{costo unitario de produccion.}\\
\bottomrule
\end{array}
\]
Las utilidades son:
\[
\Pi_s=\text{utilidad del proveedor},
\qquad
\Pi_r=\text{utilidad del minorista},
\qquad
\Pi_{SC}=\Pi_s+\Pi_r.
\]
El superindice \(D\) indica escenario descentralizado y el superindice \(C\) indica cadena centralizada.

La demanda se puede escribir como:
\[
q=100-p
\quad\Longleftrightarrow\quad
p=100-q.
\]
Por lo tanto:
\[
a=100,\qquad b=1,\qquad c=1.
\]
Aqui \(a\) es el intercepto de la demanda inversa \(p(q)=a-bq\), y \(b\) es la pendiente absoluta de esa demanda.

En la cadena descentralizada:
\[
w_D=\frac{a+c}{2}
=\frac{100+1}{2}
=50.5.
\]
La cantidad y el precio final son:
\[
q_D=\frac{a-c}{4b}
=\frac{100-1}{4}
=24.75,
\]
\[
p_D=\frac{3a+c}{4}
=\frac{3(100)+1}{4}
=75.25.
\]

Utilidad proveedor:
\[
\Pi_s^D=(w_D-c)q_D
=(50.5-1)(24.75)
=1225.125.
\]
Utilidad minorista:
\[
\Pi_r^D=(p_D-w_D)q_D
=(75.25-50.5)(24.75)
=612.5625.
\]
Utilidad total:
\[
\Pi_{SC}^D=1225.125+612.5625
=1837.6875.
\]

En la cadena centralizada:
\[
q_C=\frac{a-c}{2b}
=\frac{100-1}{2}
=49.5,
\]
\[
p_C=\frac{a+c}{2}
=\frac{100+1}{2}
=50.5.
\]
La utilidad total centralizada es:
\[
\Pi_{SC}^C=(p_C-c)q_C
=(50.5-1)(49.5)
=2450.25.
\]

Comparando:
\[
q_C=49.5>24.75=q_D,
\qquad
p_C=50.5<75.25=p_D,
\]
\[
\Pi_{SC}^C=2450.25>1837.6875=\Pi_{SC}^D.
\]

La cadena descentralizada vende menos y cobra mas caro porque proveedor y minorista agregan margen por separado. Esa perdida se llama \textbf{doble marginalizacion}.
\end{solucion}
